% SEGMENT OLP-0661-S01
% Part: proof-theory
% Chapter: cut-elimination
% Section: midsequent

\documentclass[../../../include/open-logic-section]{subfiles}

\begin{document}

\olfileid{pt}{cut}{mid}

\olsection{நடுத்தொடரணித் தேற்றமும் ஹெர்பிராண்டின் தேற்றமும்}

வெட்டு நீக்கத் தேற்றத்தின் முக்கிய விளைவுகளில் ஒன்று
நடுத்தொடரணித் தேற்றம். $\Gamma \Sequent \Delta$ இன்
!!a{proof}~$\pi$ ஒன்றில் ஒவ்வொரு !!{formula} உம் முன்னிலை
அளவையடை வடிவில் இருந்தால், $S \ident \Pi \Sequent
\Lambda$ என்ற தொடரணியைக் கொண்ட !!a{proof}~$\pi'$ ஒன்று
உள்ளது என இது கூறுகிறது. $S$ ஐ \emph{நடுத்தொடரணி} என்போம்;
$\pi'$ இல் $S$ க்கு மேலுள்ள அனுமானங்கள் அனைத்தும்
கூற்றுக்கணிப்பு அனுமானங்கள், $S$ க்குக் கீழுள்ளவை அனைத்தும்
அளவையடை அனுமானங்கள். (ஒரு !!^a{formula}~$!A$,
\[\mathsf{Q}_1x_1\dots\mathsf{Q}_nx_n\,!B(x_1,\dots,x_n)\]
என்ற வடிவில் இருந்தால் அது \emph{முன்னிலை அளவையடை வடிவில்}
உள்ளது; இதையே \emph{முன்னிலை வடிவம்} என்றும் கூறுவர்.
இங்கு ஒவ்வொரு $\mathsf{Q}_i$ உம் $\lforall$ அல்லது
$\lexists$ ஆகும்.) நடுத்தொடரணித் தேற்றத்தின் மற்றொரு
முக்கிய விளைவு ஹெர்பிராண்டின் தேற்றம். அதன் பொதுவான வடிவை
விவரிப்பது சற்றுக் கடினம்; மிகவும் பொதுவாக, !!{provable}
ஆன முதல்வரிசைத் தருக்க !!{sentence}~$!A$ ஒவ்வொன்றுக்கும்,
$!A$ ஐ !!{prove} செய்ய உதவும் கூற்றுக்கணிப்பு
மெய்மம்~$!A'$ ஒன்று உள்ளது என்பதே அதன் வடிவம்.
$!B$ அளவையடைகள் அற்றதாக இருக்கும்போது
$\lexists[x][!B(x)]$ வடிவிலான !!{sentence}s க்கான
தேற்றம் இதோ:

% SEGMENT OLP-0661-S02
\begin{thm}[ஹெர்பிராண்டின் தேற்றம்]
$!B(x)$ அளவையடைகள் அற்றது என்றும்
$\Proves \lexists[x][!B(x)]$ என்றும் கொள்க.
அப்போது $\Proves !B(t_1) \lor \dots \lor !B(t_n)$
ஆகுமாறு $t_1$, \dots, $t_n$ என்ற பதங்கள் உள்ளன.
\end{thm}

$!B(t_1) \lor \dots \lor !B(t_n)$ என்ற !!{formula} ஐ,
$\lexists[x][!B(x)]$ இன் \emph{ஹெர்பிராண்ட் பிரிப்பிணைவு}
என்போம். $!B$ இல் அளவையடைகள் இல்லாததால் இப்பிரிப்பிணைவிலும்
அவை இல்லை. ஆகவே இது !!{provable} எனில், வெட்டு இல்லாத
நிறுவல் மூலமும் !!{provable} ஆகும்; அதன் விளைவாக
கூற்றுக்கணிப்பு அனுமானங்களை மட்டுமே பயன்படுத்தும்
!!a{proof} உம் உண்டு. எனவே இது ஒரு மெய்மம்.

% SEGMENT OLP-0661-S03
ஒவ்வொரு !!{provable} !!{formula}~$\lexists[x][!B(x)]$
க்கும் ஹெர்பிராண்ட் பிரிப்பிணைவு ஒன்று இருப்பதைக் காட்டுவதே
ஹெர்பிராண்டின் தேற்றத்தின் கடினமான பகுதி. மறுதிசை எளிது:
$\lexists[x][!B(x)]$ க்கு ஹெர்பிராண்ட் பிரிப்பிணைவு
ஒன்று இருந்தால் அது நிறுவத்தக்கது. உண்மையில், அந்தப்
பிரிப்பிணைவிலிருந்து $\lexists[x][!B(x)]$ பின்தொடர்கிறது.
$n=2$ நேர்வுக்கு~$\Log{G3c}$ இல் !!a{proof} இதோ:
\[
\Axiom$!B(t_1) \fCenter \lexists[x][!B(x)], !B(t_1), !B(t_2)$
\Axiom$!B(t_2) \fCenter \lexists[x][!B(x)], !B(t_1), !B(t_2)$
\RightLabel{\LeftR{\lor}}
\BinaryInf$!B(t_1) \lor !B(t_2) \fCenter \lexists[x][!B(x)], !B(t_1), !B(t_2)$
\RightLabel{\RightR{\lexists}}
\UnaryInf$!B(t_1) \lor !B(t_2) \fCenter \lexists[x][!B(x)], !B(t_2)$
\RightLabel{\RightR{\lexists}}
\UnaryInf$!B(t_1) \lor !B(t_2) \fCenter \lexists[x][!B(x)]$
\DisplayProof
\]

% SEGMENT OLP-0661-S04
$\Sequent \lexists[x][!B(x)]$ இன் !!a{proof}~$\pi$
இலிருந்து ஹெர்பிராண்ட் பிரிப்பிணைவைப் பெறுவது எப்படி
என்பதைப் பார்ப்போம். வெட்டு இல்லாத !!{proof}~$\pi$ இல்
எல்லா $\RightR{\lexists}$ அனுமானங்களும் இறுதியில்
ஒன்றாக வருவதாக முதலில் கொள்க:
\[
\AxiomC{}
\RightLabel{$\pi_1$}
\Deduce$ \fCenter \lexists[x][!B(x)], !B(t_1), \dots, !B(t_n)$
\RightLabel{\RightR{\lexists}}
\UnaryInf$ \fCenter \lexists[x][!B(x)], !B(t_2), \dots, !B(t_n)$
\RightLabel{$(n-1) \times \RightR{\lexists}$}
\Deduce$ \fCenter \lexists[x][!B(x)]$
\DisplayProof
\]
$\lexists[x][!B(x)]$ செயல்படு வாய்பாடாக வரும்
$\pi$ இன் $\RightR{\lexists}$ அனுமானங்கள் இவையே
எனில், $\pi_1$ இல் வேறு அளவையடை அனுமானங்கள்
இருக்காது. ஏனெனில் $!B(x)$ இல் வேறு அளவையடைகள்
இல்லை; இறுதித் தொடரணியின் இடப்புறத்திலும் எந்த
!!{formula}s உம் இல்லை. வேறு சொற்களில்,
$ \fCenter \lexists[x][!B(x)], !B(t_1),
\dots, !B(t_n)$ என்பது~$\pi$ இன் நடுத்தொடரணி.
மேலும், நடுத்தொடரணிக்கு மேல் அளவையடை அனுமானங்கள்
இல்லாததால், அதன் மேலுள்ள எல்லாத் தொடரணிகளிலும்
$\lexists[x][!B(x)]$ வரும்; ஆனால் அது செயல்படு
வாய்பாடாகவோ முதன்மை வாய்பாடாகவோ இருக்காது.
எனவே~$\pi_1$ என்ற !!{proof} இலிருந்து அதை நீக்கி,
$\fCenter !B(t_1), \dots, !B(t_n)$ இன்
!!a{proof} பெறலாம். அதன் பின் $\RightR{\lor}$
ஐ $n-1$ முறை பயன்படுத்தி,
$\Sequent !B(t_1) \lor \dots \lor !B(t_n)$
என்ற ஹெர்பிராண்ட் பிரிப்பிணைவின் !!a{proof} பெறலாம்.

% SEGMENT OLP-0661-S05
ஆனால், $\Sequent \lexists[x][!B(x)]$ இன் வெட்டு
இல்லாத !!{proof} ஒன்றில்கூட எல்லா
\RightR{\lexists} அனுமானங்களும் இறுதியில் ஒரே
தொகுதியாக வரும் என்று கருத முடியாது. இந்த
\RightR{\lexists} அனுமானங்களுக்கு இடையில் ஏதாவது
$!B(t_i)$ முதன்மை வாய்பாடாக இருக்கும்
அனுமானங்கள் வரலாம். இதற்குத்தான் நடுத்தொடரணித்
தேற்றம் தேவை.

\begin{thm}[நடுத்தொடரணித் தேற்றம்]\ollabel{thm:midsequent}
$\Gamma$, $\Delta$ ஆகியவற்றிலுள்ள எல்லா
!!{formula}s உம் முன்னிலை அளவையடை வடிவில் இருக்க,
$\Gamma \Sequent \Delta$ இன்
\Log{G3c}-!!{proof}~$\pi$ ஒன்று இருந்தால்,
$\Pi \Sequent \Lambda$ என்ற தொடரணியைக் கொண்ட
!!a{proof}~$\pi'$ ஒன்றும் உள்ளது. அதற்குக்
கீழுள்ள அனுமானங்கள் அனைத்தும் அளவையடை
அனுமானங்கள்; மேலுள்ளவை அனைத்தும்
கூற்றுக்கணிப்பு அனுமானங்கள்.
\end{thm}

% SEGMENT OLP-0661-S06
\begin{proof}
  $\pi$ இல் ஓர் அளவையடை அனுமானத்தின்
  \emph{வரிசை மதிப்பு} என்பதை அதன் கீழுள்ள
  கூற்றுக்கணிப்பு அனுமானங்களின் எண்ணிக்கையாக
  வரையறுக்கிறோம். $\pi$ இன் வரிசை மதிப்பான
  $o(\pi)$, அதிலுள்ள எல்லா அளவையடை
  அனுமானங்களின் வரிசை மதிப்புகளின் கூட்டுத்தொகை.
  $o(\pi) = 0$ எனில், எந்த அளவையடை
  அனுமானத்திற்குக் கீழும் கூற்றுக்கணிப்பு
  அனுமானம் இல்லை. அளவையடை அனுமானமே இல்லை
  எனில் இறுதித் தொடரணியையே நடுத்தொடரணியாகக்
  கொள்ளலாம். அவை இருந்தால், ஒவ்வொன்றுக்கும்
  ஒரு முன்கோள் மட்டுமே இருப்பதால், மிக மேலுள்ள
  அளவையடை அனுமானம் ஒன்றே இருக்கும்; அதன்
  முன்கோளே நடுத்தொடரணி.

  $o(\pi) > 0$ எனில், வரிசை மதிப்பு $>0$ உடைய
  அளவையடை அனுமானங்களில் மிகக் கீழுள்ள ஒன்றைத்
  தேர்க. அதன் மதிப்பு $>0$ என்பதால் அதற்குக்
  கீழே ஒரு கூற்றுக்கணிப்பு அனுமானம் இருக்க வேண்டும்.
  தேர்ந்தெடுத்தது மிகக் கீழுள்ள அளவையடை
  அனுமானம் என்பதால், அதன் முடிவு மற்றோர்
  அளவையடை அனுமானத்தின் முன்கோளாக இருக்க முடியாது:
  அத்தகைய கீழ் அனுமானத்திற்கும் அதன் கீழே
  கூற்றுக்கணிப்பு அனுமானம் இருக்கும். மேலுள்ள
  அளவையடை விதியும் அதன் கீழுள்ள
  கூற்றுக்கணிப்பு விதியும் எவை என்பதைப் பொறுத்து
  பல நேர்வுகளைப் பிரிக்க வேண்டும்.

  இறுதித் தொடரணியில் முன்னிலை அளவையடை வடிவிலான
  வாய்பாடுகள் மட்டுமே இருப்பதால், அளவையடை
  அனுமானத்தின் முதன்மை வாய்பாடு அடுத்த
  கூற்றுக்கணிப்பு அனுமானத்தில் செயல்படு வாய்பாடாக
  இருக்க முடியாது. அப்படியானால், அந்த
  கூற்றுக்கணிப்பு அனுமானத்தின் முதன்மை வாய்பாட்டில்
  ஓர் இணைப்பியின் ஆட்சிக்குள் அளவையடை இருக்கும்.
  உள்-!!{formula} பண்பின்படி, அத்தகைய வாய்பாடு
  இறுதித் தொடரணியிலுள்ள !!a{formula} வின்
  உள்-!!{formula} ஆக வேண்டியிருக்கும்.
  ஆகவே அளவையடை அனுமானத்தின் முதன்மை
  !!{formula}, கீழுள்ள கூற்றுக்கணிப்பு
  அனுமானத்தின் சூழலில் ஓர் !!a{element} ஆகும்.

  இரண்டு எடுத்துக்காட்டுகளைப் பார்ப்போம்:

% SEGMENT OLP-0661-S07
முதலில், ஒரு $\RightR{\lforall}$ அனுமானம்
$\LeftR{\lif}$ விதியின் வலது முன்கோளில்
முடிகிறது என்க. அப்போது~$\pi$ இன் தொடர்புடைய
உள்-!!{proof}~$\pi_1$ பின்வருமாறு முடிகிறது:
\[
\AxiomC{}
\RightLabel{$\pi_2$}
\Deduce$\Gamma' \fCenter \Delta', \lforall[x][!A(x)], !B$
\AxiomC{}
\RightLabel{$\pi_3$}
\Deduce$!C, \Gamma' \fCenter \Delta', !A(c)$
\RightLabel{\RightR{\lforall}}
\UnaryInf$!C, \Gamma' \fCenter \Delta', \lforall[x][!A(x)]$
\RightLabel{\LeftR{\lif}}
\BinaryInf$!B \lif !C, \Gamma' \fCenter \Delta', \lforall[x][!A(x)]$
\DisplayProof
\]
$\pi$ ஒழுங்கானது என்று கொள்ளலாம். ஆகவே
தனிமாறி~$c$, $\pi_3$ க்கு வெளியே வராது;
குறிப்பாக, $!B \lif !C$ இலோ $\pi_2$ இலோ
வராது. இப்போது பின்வரும் !!{proof}~$\pi_1'$
ஐக் கருதுக:
\[
\AxiomC{}
\RightLabel{$\pi_2'$}
\Deduce$\Gamma' \fCenter \Delta', \lforall[x][!A(x)], !A(c), !B$
\AxiomC{}
\RightLabel{$\pi_3$}
\Deduce$!C, \Gamma' \fCenter \Delta', \lforall[x][!A(x)], !A(c)$
\RightLabel{\LeftR{\lif}}
\BinaryInf$!B \lif !C, \Gamma' \fCenter \Delta', \lforall[x][!A(x)], !A(c)$
\RightLabel{\RightR{\lforall}}
\UnaryInf$!B \lif !C, \Gamma' \fCenter \Delta', \lforall[x][!A(x)], \lforall[x][!A(x)]$
\DisplayProof
\]

இங்கு~$\pi_2$ இன் வலப்புறத்தில் $!A(c)$ ஐ
வலுவிழக்கச் சேர்த்துப்~$\pi_2'$ பெறப்படுகிறது.
(இதனால் புதிய அனுமானங்கள் சேராது; வரிசை
மதிப்பைப் பாதிக்கக்கூடிய அளவையடை அனுமானங்களும்
சேராது. மேலும்~$\pi_2$ இல் உள்ள தனிமாறி
நிபந்தனைகள் மீறப்படாது; ஏனெனில் $!A(c)$
இலுள்ள எந்த மாறிலியும்~$\pi_2$ இல்
தனிமாறியாகப் பயன்படுத்தப்படவில்லை.)
$c$, $!B \lif !C$ இல் வராததால்,
\RightR{\lforall} இன் தனிமாறி நிபந்தனையும்
தொடர்ந்து நிறைவேறுகிறது.

% SEGMENT OLP-0661-S08
சுருக்கத்தின் ஏற்கத்தக்கதன்மையைப் பயன்படுத்தி,
$!B \lif !C, \Gamma' \fCenter \Delta',
\lforall[x][!A(x)]$ இன்
!!a{proof}~$\pi_1''$ ஐப் பெறுகிறோம்.
$\lforall[x][!A(x)]$ க்கான
$\RightR{\Contraction}$ ஏற்கத்தக்கது என்ற
நிறுவலும், அதில் பயன்படுத்திய
$\RightR{\lforall}$ மீளாக்கத்தன்மையின்
நிறுவலும், எந்த அனுமானத்தின் வரிசை
மதிப்பையும் மாற்றவில்லை: அதிகபட்சம் சில
அனுமானங்களை மட்டுமே நீக்கின. எனவே
$\pi_1''$ இல்~$\pi_1$ ஐவிட அதிகமான
அளவையடை அனுமானங்கள் இருக்காது; மேலும்
$\pi_1''$ இன் ஒவ்வோர் அளவையடை அனுமானத்திற்கும்
கீழுள்ள கூற்றுக்கணிப்பு அனுமானங்களின் எண்ணிக்கை,
$\pi_1$ இல் அதற்குரிய அனுமானத்தின் கீழுள்ள
எண்ணிக்கைக்குச் சமம். இறுதி
\RightR{\lforall} மட்டும் விதிவிலக்கு:
$\pi$ இல் அதன் கீழிருந்த~\LeftR{\lif}
ஐ~$\pi_1''$ இல் அதற்கு மேலே நகர்த்தியுள்ளோம்.
$\pi$ இல்~$\pi_1$ க்குப் பதிலாக~$\pi_1''$
ஐ இடுக. பெறப்படும் !!{proof} இன் வரிசை
மதிப்பு குறையும்; ஏனெனில் குறைந்தபட்சம்
அந்த $\RightR{\lforall}$ அனுமானத்தின்
வரிசை மதிப்பு~$1$ ஆல் குறைந்துள்ளது.
இப்போது தொகுத்தறிதல் கருதுகோளைப் பயன்படுத்துக.

% SEGMENT OLP-0661-S09
அளவையடை அனுமானம்~\RightR{\lexists}
ஆகும் நேர்வுகள் இன்னும் எளிது: சுருக்கம்
தேவையில்லை; தனிமாறி நிபந்தனையையும்
கவனிக்க வேண்டியதில்லை. எடுத்துக்காட்டாக,
$\RightR{\lexists}$ க்குப் பின்
$\RightR{\land}$ வருகிறது என்றும், அது
இடது முன்கோளில் இருக்கிறது என்றும் கொள்க.
தொடர்புடைய உள்-!!{proof}~$\pi_1$:
\[
\AxiomC{}
\RightLabel{$\pi_2$}
\Deduce$\Gamma' \fCenter \Delta', \lexists[x][!A(x)], !A(t), !B$
\RightLabel{\RightR{\lexists}}
\UnaryInf$\Gamma' \fCenter \Delta', \lexists[x][!A(x)], !B$
\AxiomC{}
\RightLabel{$\pi_3$}
\Deduce$\Gamma' \fCenter \Delta', \lexists[x][!A(x)], !C$
\RightLabel{\RightR{\land}}
\BinaryInf$\Gamma' \fCenter \Delta', \lexists[x][!A(x)], !B \land !C$
\DisplayProof
\]
$\pi$ இல்~$\pi_1$ க்குப் பதிலாகப்
பின்வரும்~$\pi_1'$ ஐ இடுக:
\[
\AxiomC{}
\RightLabel{$\pi_2$}
\Deduce$\Gamma' \fCenter \Delta', \lexists[x][!A(x)], !A(t), !B$
\AxiomC{}
\RightLabel{$\pi_3'$}
\Deduce$\Gamma' \fCenter \Delta', \lexists[x][!A(x)], !A(t), !C$
\RightLabel{\RightR{\land}}
\BinaryInf$\Gamma' \fCenter \Delta', \lexists[x][!A(x)], !A(t), !B \land !C$
\RightLabel{\RightR{\lexists}}
\UnaryInf$\Gamma' \fCenter \Delta', \lexists[x][!A(x)], !B \land !C$
\DisplayProof
\]
இங்கு~$\pi_3$ இன் வலப்புறத்தில் $!A(t)$ ஐ
வலுவிழக்கச் சேர்த்துப்~$\pi_3'$ பெறுகிறோம்.
இந்த வலுவிழக்கச் சேர்த்தலில்~$\pi_3$ இன்
ஒவ்வொரு தொடரணியின் வலப்புறத்திலும் $!A(t)$
சேர்க்கப்படுவது மட்டுமே நடக்கிறது; எனவே
எந்த அளவையடை அனுமானத்தின் வரிசை மதிப்பும்
மாறாது. $\pi$ இலும்~$\pi'$ இலும் உள்ள
அளவையடை அனுமானங்களின் வரிசை மதிப்புகள்
ஒன்றே; ஆனால்~\RightR{\land} உடன்
இடமாற்றப்பட்ட~\RightR{\lexists}
அனுமானத்தின் மதிப்பு மட்டும்~$1$ ஆல்
குறைந்துள்ளது. தொகுத்தறிதல் கருதுகோள்
பொருந்துகிறது.
\end{proof}

% SEGMENT OLP-0661-S10
\begin{prob}
\olref{thm:midsequent} இன் நிறுவலில்,
\LeftR{\lexists} ஆனது~\RightR{\lif} இன்
முன்கோளில் முடியும் நேர்வையும்,
\LeftR{\lforall} ஆனது~\LeftR{\lor} இன்
இடது முன்கோளில் முடியும் நேர்வையும் விளக்குக.
\end{prob}

% SEGMENT OLP-0661-S11
\begin{proof}[ஹெர்பிராண்டின் தேற்றத்தின் நிறுவல்]
$\pi$ ஆனது~$\Sequent \lexists[x][!A(x)]$
இல் முடிகிறது என்க.
\olref{thm:midsequent} மூலம்~$\pi$ ஐ,
நடுத்தொடரணி~$S$ உடைய !!a{proof}~$\pi'$
ஆக மாற்றலாம். அந்த நடுத்தொடரணி
\[\Sequent \lexists[x][!A(x)], !A(t_1), \dots, !A(t_n).\]
என்ற வடிவில்தான் இருக்கும். $S$ க்கு
மேலுள்ள எல்லா அனுமானங்களும் கூற்றுக்கணிப்பு
அனுமானங்கள் என்பதால், $\lexists[x][!A(x)]$
அங்கு எந்த அனுமானத்திலும் முதன்மை
வாய்பாடாக இருக்க முடியாது. அது அணு
வாய்பாடு அல்ல; எனவே அடிக்கோளிலும்
முதன்மை வாய்பாடாக முடியாது. மேலும்
$\lexists[x][!A(x)]$, $!A(t_1)$ இன்
சரியான உள்-வாய்பாடாக முடியாது;
$!A(x)$ இல் அளவையடைகள் இல்லை என்பதை
நினைவுகொள்க. ஆகவே $S$ க்கு மேல்
செயல்படு வாய்பாடாகவும் அது வராது.
$S$ இல் முடியும் உள்-!!{proof} இலிருந்து
$\lexists[x][!A(x)]$ ஐ நீக்கினால்
$\Sequent !A(t_1), \dots, !A(t_n)$
இன் சரியான !!{proof} கிடைக்கும்.
அதிலிருந்து $\RightR{\lor}$ ஐ $n-1$
முறை பயன்படுத்தி,
$\Sequent !A(t_1) \lor \dots \lor !A(t_n)$
இன் !!a{proof} பெறலாம்.
\end{proof}

% \begin{prob}
% $\lexists[x][\lforall[y][(!A(x) \lif
% !A(y))]]$ இன் வெட்டு இல்லாத !!{proof} ஒன்றை
% $\Sequent
% \lexists[x][\lforall[y][(!A(x) \lif !A(y))]]$
% க்குக் கண்டுபிடித்து, \olref{thm:midsequent} இன்
% மாற்றத்தைப் பயன்படுத்தி நடுத்தொடரணி ஒன்றைப் பெற்று
% (தேவையெனில்), அதன் ஹெர்பிராண்ட் பிரிப்பிணைவைக் காண்க.
% \end{prob}

\end{document}
